% ============================================================
% SECCIÓN 3: Protocolo Demeter
% ============================================================
\chapter{Protocolo de Comunicación Demeter}
\label{sec:protocolo}

\section{Propósito y Filosofía de Diseño}

El Protocolo Demeter es un protocolo de comunicación binario diseñado específicamente para
la red de campo del sistema. Los protocolos estándar de texto (JSON, CSV) son demasiado
pesados para microcontroladores con memoria limitada y canales de radio estrecho. Un paquete
JSON con temperatura y humedad puede ocupar entre 40 y 80 bytes. La misma información en
el Protocolo Demeter ocupa exactamente \textbf{10 bytes}.

\begin{designnote}
Las decisiones de diseño que guían el protocolo son: \textbf{compacidad} (cada bit cuenta),
\textbf{determinismo} (la longitud de cada trama es conocida antes de recibirla),
\textbf{robustez} (verificación de integridad en toda trama) y
\textbf{extensibilidad} (hasta 254 tipos de mensaje sin modificar la estructura base).
\end{designnote}

% ─────────────────────────────────────────────────────────────
\section{Modelos Pydantic: Bridging entre Binario y JSON}
\label{sec:pydantic}

Una aportación central de la arquitectura Demeter es que el protocolo no es exclusivamente
binario: las mismas definiciones de mensaje coexisten en \textbf{dos representaciones},
gestionadas mediante \textbf{modelos Pydantic} validados y tipados:

\begin{enumerate}
    \item \textbf{Representación JSON}: usada en la comunicación entre el Frontend y el
          Backend por HTTP y WebSocket. Los campos viajan con sus nombres completos
          (e.g. \texttt{"type"}, \texttt{"target\_id"}, \texttt{"pin"}, \texttt{"value"}).
    \item \textbf{Representación Binaria}: usada en la comunicación de red de campo
          (UART entre Raspberry Pi y Gateway, y ESP-NOW entre nodos). Solo los valores
          numéricos viajan en la trama; los nombres de campo \textit{no existen en el paquete},
          reduciéndolo a su mínima expresión.
\end{enumerate}

\subsection{Diseño de Modelos Discriminados}

Cada tipo de comando del protocolo tiene su propio modelo Pydantic con un campo
\texttt{type} que actúa como \textbf{discriminador de unión}. Esto permite que al
deserializar un JSON, Pydantic identifique automáticamente de qué tipo de mensaje
se trata y lo valide con el esquema apropiado:

\begin{lstlisting}[style=trama, caption={Estructura conceptual de un modelo discriminado}]
{
  "type": "set_gpio",   <-- discriminador: identifica el subtipo exacto
  "target_id": 3,       <-- nodo destino
  "pin": 4,             <-- pin a controlar
  "value": 1,           <-- 1=ON, 0=OFF
  "duration_ms": 15000  <-- tiempo de activacion (ms)
}
\end{lstlisting}

La clave del diseño es que el campo \texttt{type} es unívoco por subtipo. Cuando el
Backend recibe este JSON (desde el Frontend o la API), instancia mediante
\texttt{TypeAdapter[AnyDemeterCommand]} el modelo correcto sin lógica condicional manual.

\subsection{Auto-descubrimiento y Clasificación Automática}

El uso del patrón de \textbf{unión discriminada} con Pydantic permite que cualquier capa
del sistema clasifique un mensaje desconocido sin necesidad de \textit{switch-case} explícitos:
la librería inspecciona el campo \texttt{type} y devuelve la instancia del modelo correcto,
ya validada y tipada. Esto tiene tres consecuencias importantes:

\begin{itemize}[leftmargin=*, label=--]
    \item \textbf{Validación zero-cost}: si un campo falta o tiene tipo incorrecto,
          Pydantic lanza una excepción detallada antes de que el sistema procese nada.
    \item \textbf{Serialización simétrica}: \texttt{model.model\_dump\_json()} produce el
          JSON que el Frontend espera; \texttt{model.to\_binary()} produce la trama que
          el ESP32 espera. La lógica de transformación reside en el modelo, no en el
          controlador.
    \item \textbf{Reutilización del esquema}: el mismo \texttt{Common/schemas.py} es
          importado por el Backend (Python) y puede referenciarse para generar SDKs cliente,
          documentación automática (OpenAPI) o tests de integración.
\end{itemize}

\subsection{Conversión Binaria: Solo los Valores}

La diferencia fundamental entre JSON y binario es que en la trama física
\textbf{los nombres de los campos no existen}. El orden y tamaño de cada byte es implícito
al tipo de comando (definido por el \texttt{CMD\_ID} de la cabecera). Por ejemplo, para un
\texttt{SET\_GPIO} el payload es exactamente 3 bytes: \texttt{[PIN, VALUE, FLAGS]}.

La serialización binaria rellena cada byte en orden, con el tamaño y signo definidos por
el esquema del protocolo (enteros sin signo, punt fijo para float, Little Endian).

% ─────────────────────────────────────────────────────────────
\section{Comunicación HTTP y WebSocket en el Sistema}
\label{sec:http_ws}

\subsection{Flujo de un Comando desde el Navegador}

Cuando un operador interactúa con el dashboard activo, los comandos siguen un flujo
completamente desacoplado de la red de campo:

\begin{enumerate}
    \item El \textbf{Frontend React} compone un objeto JSON conforme al esquema del modelo
          Pydantic correspondiente (e.g. \texttt{set\_gpio}) y lo envía mediante una petición
          \textbf{HTTP POST} a \texttt{/api/command}.
    \item El \textbf{Backend FastAPI} recibe el JSON, lo valida automáticamente con el
          \texttt{TypeAdapter} discriminado, y publica el comando en el canal Redis
          \texttt{demeter:commands}. Devuelve \texttt{202 Accepted} inmediatamente.
    \item El \textbf{WS Dispatcher} (tarea asyncio en el Backend) consume el mensaje de Redis
          y lo retransmite por WebSocket a la Raspberry Pi, identificada como
          \texttt{raspberry\_gateway}.
    \item La \textbf{Raspberry Pi} recibe el JSON por WebSocket, lo convierte a trama binaria
          mediante el mismo modelo Pydantic en Python, y lo envía por UART al Gateway ESP32.
    \item El \textbf{Gateway ESP32} reenvía la trama binaria por ESP-NOW al nodo actuador
          destino, que ejecuta la acción y confirma con \texttt{CMD\_ACK}.
\end{enumerate}

\subsection{Flujo de Telemetría en Tiempo Real}

El camino inverso (del sensor al dashboard) funciona de forma reactiva:

\begin{enumerate}
    \item El \textbf{Nodo Sensor} emite tramas binarias por ESP-NOW en el intervalo configurado.
    \item El \textbf{Gateway ESP32} las reenvía por UART a la Raspberry Pi.
    \item La \textbf{Raspberry Pi} las decodifica (Protocolo Demeter en Python), construye un
          JSON y lo publica por WebSocket hacia el Backend.
    \item El \textbf{Backend} almacena la medición en Redis (caché 60~s) y en PostgreSQL
          (persistencia), y hace \textit{broadcast} por WebSocket a todos los clientes
          Frontend suscritos.
    \item El \textbf{Frontend} recibe el push y actualiza los gráficos sin ninguna petición
          adicional por parte del usuario.
\end{enumerate}

\begin{designnote}
\textbf{Separación de canales}: Los comandos viajan por \textbf{HTTP POST} (fiables, con
respuesta de estado). La telemetría viaja por \textbf{WebSocket} (baja latencia, push).
Para la Raspberry Pi, el canal WebSocket es exclusivo y bidireccional. Para el Frontend,
el WebSocket es de solo lectura; sus comandos siempre van por HTTP.
\end{designnote}

% ─────────────────────────────────────────────────────────────
\section{Estructura de una Trama}

Toda comunicación sigue el mismo formato de cabecera fija seguida de un payload variable:

\begin{table}[H]
    \centering
    \caption{Campos de la cabecera Demeter}
    \begin{tabular}{clp{8cm}}
        \toprule
        \textbf{Campo}  & \textbf{Tamaño} & \textbf{Descripción} \\
        \midrule
        SYNC    & 1 byte & Siempre \texttt{0xAA}. Marca el inicio de cualquier trama válida. \\
        FLAGS   & 1 byte & Reservado. Actualmente \texttt{0x00}. \\
        SRC\_ID & 1 byte & Identificador numérico del nodo emisor (1--253). \\
        DST\_ID & 1 byte & Destino: \texttt{0x00}=servidor, \texttt{0xFF}=difusión. \\
        CMD\_ID & 1 byte & Tipo de mensaje (ver catálogo). \\
        LENGTH  & 1 byte & Número de bytes del payload (0--255). \\
        PAYLOAD & N bytes & Datos específicos del tipo de mensaje, formato Little Endian. \\
        CRC     & 1 byte & Suma módulo 256 de FLAGS hasta fin del payload. \\
        \bottomrule
    \end{tabular}
    \label{tab:cabecera_demeter}
\end{table}

\textbf{Tamaños:} mínimo 7 bytes (sin payload) · máximo teórico 262 bytes · típico 10--14 bytes.

% ─────────────────────────────────────────────────────────────
\section{Catálogo de Mensajes}

\begin{table}[H]
    \centering
    \caption{Catálogo completo de comandos del Protocolo Demeter v2}
    \begin{tabular}{clcp{5cm}}
        \toprule
        \textbf{CMD\_ID} & \textbf{Nombre} & \textbf{Payload (bytes)} & \textbf{Descripción} \\
        \midrule
        \texttt{0x01} & PING            & 0  & Keep-alive. Espera ACK. \\
        \texttt{0x02} & ACK             & 0--1 & Confirmación con contexto de sesión opcional. \\
        \texttt{0x03} & NACK            & 1  & Rechazo con código de error. \\
        \texttt{0x04} & SYN             & 0  & Inicio del handshake (paso 1/3). \\
        \texttt{0x05} & SYN-ACK         & 0  & Confirmación de disponibilidad (paso 2/3). \\
        \texttt{0x0A} & ROUTE\_ADD      & 7  & Registra una nueva ruta en la tabla del Gateway. \\
        \texttt{0x0B} & TEMP\_HUM\_REPORT & 4 & Datos de temperatura y humedad (punto fijo ×100). \\
        \texttt{0x0C} & PIN\_REPORT     & 2  & Estado de un pin tras ejecución de un comando. \\
        \texttt{0x0D} & SYSTEM\_REPORT  & 8  & Modo de operación y nivel de batería en mV. \\
        \texttt{0x10} & SET\_GPIO       & 3  & Orden de cambio de estado de un pin. \\
        \texttt{0x11} & SET\_PWM        & 3  & Señal analógica (control de velocidad/intensidad). \\
        \texttt{0x20} & GET\_SENSORS    & 0  & Solicitud de lectura inmediata de sensores. \\
        \texttt{0x30} & EXEC\_SEQUENCE  & Variable & Lista de pasos de actuación con temporización. \\
        \bottomrule
    \end{tabular}
    \label{tab:catalogo_comandos}
\end{table}

% ─────────────────────────────────────────────────────────────
\section{Codificación de Datos}

\subsection{Temperatura y Humedad (Punto Fijo ×100)}

Los valores se multiplican por 100 y se codifican como enteros con signo de 16 bits
en formato Little Endian:

\[
T = 24{,}57\,^{\circ}\text{C} \rightarrow 2457 = \texttt{0x0999} \rightarrow \text{bytes: } \texttt{[99\;09]}
\]
\[
H = 63{,}20\,\% \rightarrow 6320 = \texttt{0x18B0} \rightarrow \text{bytes: } \texttt{[B0\;18]}
\]

\subsection{Duraciones en Secuencias (32 bits Little Endian)}

Las duraciones se expresan en milisegundos como entero de 32 bits sin signo:
\[
2500\,\text{ms} = \texttt{0x000009C4} \rightarrow \text{bytes: } \texttt{[C4\;09\;00\;00]}
\]

% ─────────────────────────────────────────────────────────────
\section{Ejemplos de Tramas Reales}

\subsection{Telemetría de Sensor (TEMP\_HUM\_REPORT)}
Nodo 2 envía T=24,57°C y H=63,20\% al servidor:

\begin{lstlisting}[style=trama, caption={Trama TEMP\_HUM\_REPORT completa}]
Pos: [0]   [1]   [2]   [3]   [4]   [5]   [6]   [7]   [8]   [9]
Hex: AA     00    02    00    0B    04    99    09    B0    18   CRC
     SYNC  FLAGS  SRC   DST   CMD   LEN  T_Lo  T_Hi  H_Lo  H_Hi
                  (2)   (0) TH_RPT  (4)  ---- 24.57 C ----  -63.20%-
\end{lstlisting}

\subsection{Comando de Actuación (SET\_GPIO)}
Servidor ordena al nodo 3 encender el pin 4 (bomba):

\begin{lstlisting}[style=trama, caption={Trama SET\_GPIO}]
Pos: [0]   [1]   [2]   [3]   [4]   [5]   [6]   [7]   [8]   [9]
Hex: AA     00    00    03    10    03    04    01    00    CRC
     SYNC  FLAGS  SRC   DST   CMD   LEN   PIN  VAL  FLAGS
                  (0)   (3) SET_IO  (3)   (4) HIGH    (0)
\end{lstlisting}

\subsection{Inicio de Handshake (SYN)}
Nodo sensor 2 solicita sesión con el Gateway (nodo 1):

\begin{lstlisting}[style=trama, caption={Trama SYN mínima --- 7 bytes}]
Pos: [0]   [1]   [2]   [3]   [4]   [5]   [6]
Hex: AA     00    02    01    04    00    CRC
     SYNC  FLAGS  SRC   DST   CMD   LEN
                  (2)   (1)   SYN   (0)
\end{lstlisting}

% ─────────────────────────────────────────────────────────────
\section{Algoritmos de Serialización y Deserialización}

\subsection{Serialización}

El proceso de serialización transforma los datos de aplicación en una trama binaria
completamente estructurada. Toma el comando o reporte, convierte sus campos al
formato numérico requerido (punto fijo para temperatura, Little Endian para enteros),
construye el payload en un buffer, prepende la cabecera y calcula el CRC.

\subsection{Deserialización y Validación}

El \textit{pipeline} de deserialización actúa como barrera de seguridad: verifica
progresivamente el byte de inicio (SYNC), la longitud declarada, el CRC y el destino.
Solo si todas las comprobaciones pasan se decodifica el payload y se lanza el evento
interno asociado al tipo de mensaje.

\subsection{Manejo de Errores}

\begin{warningbox}
El protocolo adopta la filosofía de \textbf{fallo silencioso}: una trama corrupta no genera
ninguna respuesta de error. Esto evita tormentas de retroalimentación ante interferencias de
radio. El NACK se reserva para errores de lógica de aplicación (pin no válido, modo IDLE activo).
\end{warningbox}

% ─────────────────────────────────────────────────────────────
\section{Handshake TCP-like}

Antes de operar, cada nodo negocia una sesión con el Gateway mediante un triple
\textit{handshake}: el emisor envía SYN, espera SYN-ACK con \textit{timeout} de 2~s
y hasta 3 reintentos, y confirma con ACK. Solo tras este acuerdo comienzan los reportes.

% ─────────────────────────────────────────────────────────────
\section{Flujo Extremo a Extremo}

Una trama parte del Nodo Sensor en formato binario por ESP-NOW, llega al Gateway,
se reenvía por UART a la Raspberry Pi, que la decodifica y publica como JSON al Backend.
El Backend la persiste en PostgreSQL y distribuye en tiempo real a todos los clientes
conectados por WebSocket, con reacción visual a la décima de segundo.

% ─────────────────────────────────────────────────────────────
\section{Consideraciones de Tiempo y Saturación}

\subsection{Tiempos de Transmisión}

\begin{table}[H]
    \centering
    \caption{Parámetros de tiempo de transmisión del protocolo}
    \begin{tabular}{lc}
        \toprule
        \textbf{Parámetro} & \textbf{Valor} \\
        \midrule
        Tiempo de transmisión por trama (10B) & $\approx$ 0,5 ms \\
        Latencia extremo a extremo (RF + serial) & 3 -- 10 ms \\
        Intervalo de reporte del sensor (configurado) & 5.000 ms \\
        Ocupación del canal por nodo sensor & < 0,01\% \\
        \bottomrule
    \end{tabular}
    \label{tab:tiempos}
\end{table}

\subsection{Ciclo Energético del Nodo Sensor}

El nodo sensor efectúa un intenso ahorro de recursos invocando Deep Sleep del ESP32 entre
reportes. Un nodo configurado con reportes cada 60~s puede operar durante más de un año
con una batería de 3000 mAh.

\begin{warningbox}
Los nodos actuadores y el Gateway deben permanecer activos de forma continua y requieren
alimentación estable. Solo el Nodo Sensor puede operar con batería durante periodos prolongados.
\end{warningbox}

\subsection{Limitaciones Conocidas}

\begin{table}[H]
    \centering
    \caption{Limitaciones conocidas y decisiones de diseño adoptadas}
    \begin{tabular}{p{4.5cm}p{4cm}p{5.5cm}}
        \toprule
        \textbf{Limitación} & \textbf{Impacto} & \textbf{Decisión de diseño} \\
        \midrule
        Sin confirmación de entrega en ESP-NOW &
        Una trama puede perderse sin aviso &
        Aceptable para telemetría; comandos críticos usan ACK de capa de aplicación. \\
        Máx. 250 bytes por trama ESP-NOW &
        Secuencias limitadas a $\approx$30 pasos por trama &
        Suficiente para los casos actuales; secuencias largas se pueden particionar. \\
        Sin cifrado nativo &
        Tramas legibles por cualquier ESP32 cercano &
        AES-128 en payload es mejora planificada para v3.0. \\
        Tabla de rutas fija (20 entradas) &
        Sin admisión dinámica sin reiniciar &
        El comando ROUTE\_ADD permite registrar rutas en caliente. \\
        \bottomrule
    \end{tabular}
    \label{tab:limitaciones}
\end{table}
